\begin{figure}[H]
    \centering
    \begin{tikzpicture}[>=Stealth, arraynode/.style={draw, minimum width=7mm, minimum height=6mm, outer sep=0, inner sep=0}]

        \node[arraynode, fill=lightgray] (a1) {000};
        \node[arraynode, fill=lightorange, right=0 of a1] (a2) {001};
        \node[arraynode, fill=lightyellow, right=0 of a2] (a3) {000};
        \node[arraynode, right=0 of a3, path picture={
        \shade[left color=lightorange, right color=lightred]
                (path picture bounding box.south west)
                rectangle (path picture bounding box.north east);
        }] (a4) {101};
        \node[arraynode, fill=lightyellow, right=0 of a4] (a5) {000};
        \node[arraynode, right=0 of a5, path picture={
        \shade[left color=lightorange, right color=lightred]
                (path picture bounding box.south west)
                rectangle (path picture bounding box.north east);
        }] (a6) {000};
        \node[arraynode, fill=lightyellow, right=0 of a6] (a7) {000};
        \node[arraynode, fill=lightorange, right=0 of a7] (a8) {000};
        \node[arraynode, fill=lightyellow, right=0 of a8, path picture={
        \shade[left color=lightred, right color=lightpink]
                (path picture bounding box.south west)
                rectangle (path picture bounding box.north east);
        }] (a9) {110};
        \node[arraynode, fill=lightpink, right=0 of a9] (a10) {000};
        \node[arraynode, fill=lightpurple, right=0 of a10] (a11) {111};
        \node[arraynode, fill=lightyellow, right=0 of a11] (a12) {000};
        \node[arraynode, right=0 of a12, path picture={
        \shade[left color=lightred, right color=lightpink]
                (path picture bounding box.south west)
                rectangle (path picture bounding box.north east);
        }] (a13) {000};
        \node[arraynode, fill=lightpink, right=0 of a13] (a14) {000};
        \node[arraynode, fill=lightpurple, right=0 of a14] (a15) {000};

        \draw[thick] 
            ([yshift=4pt]a1.north west) -- 
            ([yshift=-4pt]a1.south west);
        
        \foreach \i in {1,3,5,7,9,11,13,15} {
            \draw[thick] 
                ([yshift=4pt]a\i.north east) -- 
                ([yshift=-4pt]a\i.south east);
        }

        \node[circle, draw, above=12mm of a4, fill=lightorange] (x) {$x$};
        \draw[->] (x) to (a2.north);
        \draw[->] (x) to (a4.north);
        \draw[->] (x) to (a6.north);
        \draw[->] (x) to (a8.north);
        
        \node[circle, draw, above=12mm of a7, fill=lightred] (y) {$y$};
        \draw[->] (y) to (a4.north);
        \draw[->] (y) to (a6.north);
        \draw[->] (y) to (a9.north);
        \draw[->, dashed, lightgray] (y) to[out=150, in=110] (a1.north);

        \node[circle, draw, above=12mm of a10, fill=lightpink] (w) {$w$};
        \draw[->] (w) to (a9.north);
        \draw[->] (w) to (a10.north);
        \draw[->] (w) to (a13.north);
        \draw[->] (w) to (a14.north);

        \node[circle, draw, above=12mm of a13, fill=lightpurple] (z) {$z$};
        \draw[->] (z) to (a11.north);
        \draw[->] (z) to (a13.north);
        \draw[->] (z) to (a15.north);
        \draw[->, dashed, lightgray] (z) to[out=160, in=110] (a1.north);
        
    \end{tikzpicture}
    \caption{Non-uniform binary fuse filter with phantom entry, 14-entry array, 2-entry segments, and 3-bit fingerprints. Vertical bars mark segment boundaries.
    The phantom entry is colored gray. The filter contains keys $\mathcolorbox{lightorange}{x}$, $\mathcolorbox{lightred}{y}$, $\mathcolorbox{lightpink}{w}$ and 
    $\mathcolorbox{lightpurple}{z}$, with fingerprints $\mathcolorbox{lightorange}{100}$, $\mathcolorbox{lightred}{011}$, $\mathcolorbox{lightpink}{110}$ and 
    $\mathcolorbox{lightpurple}{111}$, respectively. Entries hashed to by keys are marked with arrows and the keys' colors. Each key hashes to three or four 
    consecutive segments. Dashed gray arrows indicate hashes to the phantom entry.}
    \label{fig:csc592_bigdata_nubff}
\end{figure}
